% !PRINTED 156
\weberpage{156}
%<<CONT
of definite dimension form an (infinite) \gesperrt{group} (\S.\,1). If we pick out of this totality any set which is closed in itself, in such a way that any two of its elements yield by composition an element of the same set, then this set likewise forms a \gesperrt{group}, which may be finite or infinite.

If $L$ denotes a fixed substitution, then from every substitution $A$ of the same dimension one can derive a substitution

\weq{21}{A' = L^{-1}AL}

which is called the \gesperrt{transform of $A$ by $L$}.

If we set

\weq{22}{(x) = L(x'), \quad (y) = L(y'),}

then it follows from (4):

\weq{23}{(y') = A'(x'),}

so that the passage to the transformed substitution is equivalent to the simultaneous transformation of both sequences of variables by $L^{-1}$.

If

\[
A' = L^{-1}AL, \quad B' = L^{-1}BL,
\]

then from the laws of composition it follows that

\[
A'B' = L^{-1}ABL,
\]

and from this therefore the theorem:

\par
7.\enspace\gesperrt{If $A$ runs through the substitutions of a group, then for fixed $L$ the transform $L^{-1}AL$ runs through the substitutions of an isomorphic group.}

The \gesperrt{transposed substitutions} must be carefully distinguished from the inverse substitutions.

Namely, from every substitution $A$ one obtains another definite substitution, called the transposed substitution to $A$, which we shall for the moment denote by $A_1$, by making the rows of $A$ into columns and conversely; that is, by interchanging in (2) the upper and lower indices of the $a$.

If in the sum (11), $\sum^i b_i^{(k)} a_h^{(i)}$, one interchanges the lower with the upper indices and at the same time $a$ with $b$, one obtains an expression which by (11) is equal to $c_k^{(h)}$.
